\chapter*{Заключение}
\addcontentsline{toc}{chapter}{Заключение}
\label{ch:conclusion}

В ходе выполнения лабораторного практикума были изучены и практически реализованы основные технологии построения, маршрутизации и защиты локальных вычислительных сетей. Были исследованы механизмы сегментации сети с использованием VLAN, предотвращения сетевых петель и повышения отказоустойчивости с помощью STP/RSTP, технологии маршрутизации и трансляции сетевых адресов, а также реализация атаки MAC-spoofing на виртуальную сетевую инфраструктуру.

В процессе работы были настроены и проверены статическая и динамическая маршрутизация с использованием OSPF, выполнен анализ сетевого трафика средствами Wireshark и исследованы методы выявления подозрительной сетевой активности. Также было реализовано защищённое взаимодействие между удалёнными локальными сетями с использованием GRE и IPSec, а работоспособность соединения подтверждена проверкой маршрутизации, доступности узлов и анализом ESP-трафика.

Поставленные задачи исследования выполнены в полном объёме. Полученные результаты демонстрируют возможность применения рассмотренных технологий для построения и защиты сетевой инфраструктуры, а полученные практические навыки могут использоваться при настройке, диагностике и администрировании корпоративных компьютерных сетей.

Выполненная работа позволила закрепить теоретические знания и получить практические навыки проектирования, настройки, диагностики и защиты локальных вычислительных сетей, а также продемонстрировала практическую применимость рассмотренных сетевых технологий.
